\documentclass{article}
\usepackage{fontspec}
\usepackage{amsmath}
\usepackage[a4paper,margin=1.5cm]{geometry}
\setmainfont{Libertinus Serif}
\setlength{\parindent}{0pt}

\begin{document}
Язык лямбда-выражений с именованными определениями.

\textbf{Пример:}
\begin{verbatim}
let S = \x y z.x z (y z)
let K = \x y.x
S K K
\end{verbatim}
\textbf{Грамматика:}
\[\begin{aligned}
Program &::= Definitions\ Expr \\
Definitions &::= Definition\ Definitions \mid Definition \\
Definition &::= \texttt{let}\ Var\ =\ Expr \\
Expr &::= Lambda \mid Application \\
Lambda &::= \texttt{\textbackslash}\ Vars\ .\ Expr \\
Vars &::= Var\ Vars \mid Var \\
Application &::= Atom\ ApplicationTail \\
ApplicationTail &::= Atom\ ApplicationTail \mid Atom \\
Atom &::= Var \mid (\ Expr\ )
\end{aligned}\]
\textbf{Свойства:}
\textbf{Аппликация:}
\[x\ y\ z = ((x\ y)\ z)\]
\textbf{Множественные параметры:}
\[\texttt{\textbackslash x y z. M}=
\texttt{\textbackslash x.(\textbackslash y.(\textbackslash z. M))}\]
\textbf{Приоритет:}
\[\text{Atom} > \text{Application} > \text{Lambda}\]
\texttt{let} не является допустимым именем переменной.

\textbf{Вывод:}

\textbf{Пример 1 (S K K):}
\[\begin{aligned}
S\ K\ K
&\Rightarrow Application \\
&\Rightarrow Atom\ ApplicationTail \\
&\Rightarrow S\ K\ ApplicationTail \\
&\Rightarrow S\ K\ K
\end{aligned}\]
\vspace{0.5em}
\textbf{Пример 2 (лямбда-выражение):}
\[\begin{aligned}
\texttt{\textbackslash x y. x y}
&\Rightarrow Lambda \\
&\Rightarrow \texttt{\textbackslash}\ x\ y\ .\ Expr \\
&\Rightarrow \texttt{\textbackslash}\ x\ y\ .\ x\ y
\end{aligned}\]
\vspace{0.5em}
\textbf{Пример 3 (let K):}
\[\begin{aligned}
\texttt{let K = \textbackslash x y.x}
&\Rightarrow Definition \\
&\Rightarrow \texttt{let}\ K\ =\ Lambda \\
&\Rightarrow \texttt{let}\ K\ =\ \texttt{\textbackslash} x\ y . x
\end{aligned}\]
\end{document}